\scanpage{127}\mspage{72}
rappelées dans $(*)$ p.~70), soit pour $W_S^g$ (quand ces catégories font partie des quatre du côté droit, rappelées dans $(*')$), lesquels \deleted{\unsure{équiva}} localiseurs coïncident encore avec $W_S^f$ sauf dans le seul cas de la catégorie $\mathrm{Cat}_{/S}$. Tous les foncteurs ainsi obtenus se réduisent, à équivalence près, à deux, suivant que la catégorie localisée but est $\mathrm{HOT}_W(S)$ (cas de \deleted{\unsure{foncteur} \illegible{1} du côté} \inserted{la localisation} $W_S^d$), ou qu'elle est $\mathrm{HOT}_W^\circ(S)$ (cas de la localisation $W_S^g$). On trouve ainsi des foncteurs (en sens inverse de $\unsure{\rho|p}$ et de $\sigma$)
\diagram{\begin{tikzcd}[ampersand replacement=\&,row sep=large]
 {}\&\mathrm{HOT}_W^{\unsure{\mathrm{lc}|\ell\mathrm{c}}}(S)\arrow[dl,"i"']\arrow[dr,"j"]\&{}\\
 \mathrm{HOT}_W(S)\&{}\&\mathrm{HOT}_W^\circ(S)\,.
\end{tikzcd}}
Les foncteurs $\deleted{\unsure{R}}\rho,\sigma$ s'interchangent par dualité $S\mapsto S^\circ$, et de même pour $i,j$.

\textcircled{IV} Les foncteurs $\rho,\deleted{\unsure{j}}i$ sont adjoints l'un de l'autre dans cet ordre, i.e. on a pour deux objets \deleted{$X,Y$ de $\mathrm{Cat}$} $\xi\in\mathrm{Ob}\,\mathrm{HOT}_W(S)$, $\unsure{\eta|y}\in\mathrm{Ob}\,\mathrm{HOT}^{\mathrm{lc}}_W(S)$
\[
 \mathrm{Hom}(\rho(\deleted{\unsure{\eta}}\xi),\deleted{\unsure{\xi}}\eta)\simeq
 \mathrm{Hom}(\deleted{\unsure{\eta}}\xi,\deleted{\unsure{j}}i(\deleted{\unsure{\xi}}\eta))\deleted{\unsure{,}}\quad\deleted{\unsure{\xi}}
\]
\deleted{d'où \unsure{il}} \deleted{\unsure{et}} (\unsure{isom|isom.} de bifoncteurs). Et dualement les foncteurs $\sigma,j$ \deleted{\unsure{,}} sont adjoints l'un de l'autre dans cet ordre\,: pour $\xi\in\mathrm{Ob}\,\mathrm{HOT}^\circ_W(S)$, $\eta\in\mathrm{Ob}\,\mathrm{HOT}^{\mathrm{lc}}_W(S)$, on a
\[
 \mathrm{Hom}(\xi,j(\eta))\simeq\mathrm{Hom}(\sigma(\xi),\eta).
\]
% archive sheet number: 126
